\begin{proposition}[异手征旋量的向量表示]
四维复化 Lorentz 代数的表示满足 $(\tfrac12,0)\otimes(0,\tfrac12)=(\tfrac12,\tfrac12)$。通过 $x^\mu\sigma_\mu$ 的二乘二矩阵实现，此四维空间就是向量表示。
\end{proposition}

\begin{proof}
\eqnrefs{eq:feqo-sl2c-chiral-dictionary} 固定左手作用矩阵 $B=(A^\dagger)^{-1}$。取两条左手旋量 $\xi,\chi$，它们的外积满足
\begin{equation*}
 Y=\xi\chi^\dagger,\qquad Y'=BYB^\dagger.
\end{equation*}
右侧共轭指标属于右手等价表示；单个外积通常只有秩一，一般双旋量应取这类外积的线性组合。任意复 $2\times2$ 矩阵都可唯一展开为
\begin{equation*}
 Y=V_\mu\sigma^\mu=V^0I_2-\mathbf V\cdot\boldsymbol\sigma_{\rm P},
 \qquad V^\mu=\frac12\Tr(\bar\sigma^\mu Y).
\end{equation*}
定义与它对应的矩阵 $X=V^0I_2+\mathbf V\cdot\boldsymbol\sigma_{\rm P}$。对二阶矩阵，伴随矩阵是线性映射
\begin{equation*}
 \operatorname{adj}X=(\Tr X)I_2-X=Y.
\end{equation*}
又因 $\det A=1$，
\begin{align*}
 \operatorname{adj}(AXA^\dagger)
 &=(A^\dagger)^{-1}(\operatorname{adj}X)A^{-1}\\
 &=BYB^\dagger.
\end{align*}
所以 $Y$ 的手征双旋量变换与 $X'=AXA^\dagger$ 诱导的是同一个四矢量变换，不能把左手矩阵 $B$ 与右手矩阵 $A$ 混写。其分量为
\begin{equation*}
 V'^\mu=\frac12\Tr(\bar\sigma^\mu B\sigma_\nu B^\dagger)V^\nu
 =\Lambda^\mu{}_{\nu}(A)V^\nu.
\end{equation*}
厄米矩阵对应实四矢量，一般复矩阵对应复化的四矢量。维数为四且上式是显式交织映射，因此
\begin{equation*}
 (\tfrac12,0)\otimes(0,\tfrac12)=(\tfrac12,\tfrac12).
\end{equation*}
行列式 $\det Y=(V^0)^2-\mathbf V^2$ 也被保持，提供独立的 Minkowski 范数核验。
\end{proof}

\begin{proposition}[同手征旋量与二形式分解]
有 $(\tfrac12,0)^{\otimes2}=(1,0)\oplus(0,0)$ 及其右手征对应式。复化反对称二阶张量为 $(1,0)\oplus(0,1)$，分别对应 Lorentz Hodge 星的两个 $\pm\ii$ 本征空间。
\end{proposition}

\begin{proof}
对两个左手 Weyl 指标，任意二阶对象 $T_{\alpha\beta}$ 都有唯一的对称--反对称分解
\begin{equation*}
 T_{\alpha\beta}=T_{(\alpha\beta)}+T_{[\alpha\beta]}.
\end{equation*}
二维旋量空间中反对称二阶张量只有一个独立分量，因此必可写成
\begin{equation*}
 T_{[\alpha\beta]}=\phi\,\epsilon_{\alpha\beta}.
\end{equation*}
为什么这个系数 $\phi$ 是 Lorentz 标量？因为 $A\in SL(2,\mathbb C)$ 满足
\begin{equation*}
 A_\alpha{}^\gamma A_\beta{}^\delta
 \epsilon_{\gamma\delta}
 =(\det A)\epsilon_{\alpha\beta}
 =\epsilon_{\alpha\beta}.
\end{equation*}
所以反对称的一维子空间保持不变，正是 $(0,0)$。剩下的对称矩阵满足
$T_{(11)},T_{(12)},T_{(22)}$ 三个独立复分量，并在
\begin{equation*}
 T_{(\alpha\beta)}\longmapsto
 A_\alpha{}^\gamma A_\beta{}^\delta
 T_{(\gamma\delta)}
\end{equation*}
下封闭，因此构成三维不可约表示 $(1,0)$。于是从实际指标分解得到
\begin{equation*}
 (\tfrac12,0)\otimes(\tfrac12,0)
 =(0,0)\oplus(1,0).
\end{equation*}
右手带点指标完全平行，给出
\begin{equation*}
 (0,\tfrac12)\otimes(0,\tfrac12)
 =(0,0)\oplus(0,1).
\end{equation*}

现在转向四维反对称张量 $F_{\mu\nu}=-F_{\nu\mu}$。它有
$\binom42=6$ 个独立分量。复化以后定义 Hodge 算符
\begin{equation*}
 (\star F)^{\mu\nu}
 =\frac12\epsilon^{\mu\nu\rho\sigma}F_{\rho\sigma}.
\end{equation*}
在当前 Lorentz 度规和 $\epsilon^{0123}=+1$ 约定下，二形式满足
$\star^2=-I$，所以其本征值是 $\pm\ii$。相应投影算符为
\begin{equation*}
 \mathcal P_\pm=\frac12(I\mp\ii\star),
 \qquad
 F_\pm\equiv\mathcal P_\pm F,
 \qquad
 \star F_\pm=\pm\ii F_\pm.
\end{equation*}
因为
$\mathcal P_+\mathcal P_-=0$、
$\mathcal P_++\mathcal P_-=I$，六维复二形式空间严格分成两个互补的三维子空间：
\begin{equation*}
 \Lambda^2_\mathbb C
 =\Lambda^2_+\oplus\Lambda^2_-.
\end{equation*}
Lorentz 变换与 Hodge 算符对易，所以两个三维子空间各自保持不变。与上面同手征双旋量的两个三维不可约表示比较，分别对应 $(1,0)$ 与 $(0,1)$；因此
\begin{equation*}
 F^{\mu\nu}\in(1,0)\oplus(0,1),
\end{equation*}
即正文\eqnrefs{eq:twoform-chiral-decomposition-r68}。对电磁场，$F_\pm$ 的三个独立复分量可等价组织成 $\mathbf E\pm\ii c\mathbf B$；这不是另外引入的场，而是同一个六分量 $F_{\mu\nu}$ 在 Hodge 本征子空间中的两种手征投影。
\end{proof}

\begin{lemma}[四个 Dirac 矩阵的迹]
若 $\{\gamma^\mu,\gamma^\nu\}=2\eta^{\mu\nu}I$ 且 $\operatorname{tr}I=N_s$，则 $\operatorname{tr}(\gamma^\mu\gamma^\nu\gamma^\rho\gamma^\sigma)=N_s(\eta^{\mu\nu}\eta^{\rho\sigma}-\eta^{\mu\rho}\eta^{\nu\sigma}+\eta^{\mu\sigma}\eta^{\nu\rho})$。四维 Dirac 表示取 $N_s=4$。
\end{lemma}

\begin{proof}
未偏振费米子振幅在完成自旋求和后会变成 Dirac 矩阵迹。正文已经把 Clifford 代数固定为\eqnrefs{eq:clifford-general}；这里仅用反对易关系与迹循环性，把它化成正文\eqnrefs{eq:trace-two-gamma-r9,eq:trace-four-gamma-r9}，从而得到后续树级截面最常用的两个迹公式。



由已经建立的关系
\begin{equation*}
 \{\gamma^\mu,\gamma^\nu\}=2\eta^{\mu\nu}\id(),
 \qquad
 \Tr\id()=4.
\end{equation*}
对 Clifford 关系取迹：
\begin{align*}
 \Tr(\gamma^\mu\gamma^\nu)
 +\Tr(\gamma^\nu\gamma^\mu)
 &=2\eta^{\mu\nu}\Tr\id(),\\
 2\Tr(\gamma^\mu\gamma^\nu)
 &=8\eta^{\mu\nu},
\end{align*}
其中第二行使用迹的循环性。故
\begin{equation*}
 {\Tr(\gamma^\mu\gamma^\nu)=4\eta^{\mu\nu}.}
 \end{equation*}
对四个 gamma，先把最左边 $\gamma^\mu$ 逐次移过后面的矩阵：
\begin{align*}
 T^{\mu\nu\rho\sigma}
 &\equiv\Tr(\gamma^\mu\gamma^\nu\gamma^\rho\gamma^\sigma),\\
 &=\Tr[(2\eta^{\mu\nu}-\gamma^\nu\gamma^\mu)
 \gamma^\rho\gamma^\sigma],\\
 &=2\eta^{\mu\nu}\Tr(\gamma^\rho\gamma^\sigma)
 -\Tr(\gamma^\nu\gamma^\mu\gamma^\rho\gamma^\sigma).
\end{align*}
再在第二项中把 $\gamma^\mu$ 移过 $\gamma^\rho$，
\begin{align*}
 T^{\mu\nu\rho\sigma}
 &=8\eta^{\mu\nu}\eta^{\rho\sigma}
 -2\eta^{\mu\rho}\Tr(\gamma^\nu\gamma^\sigma)
 +\Tr(\gamma^\nu\gamma^\rho\gamma^\mu\gamma^\sigma).
\end{align*}
最后用迹的循环性
\begin{equation*}
 \Tr(\gamma^\nu\gamma^\rho\gamma^\mu\gamma^\sigma)
 =\Tr(\gamma^\mu\gamma^\sigma\gamma^\nu\gamma^\rho)
\end{equation*}
并重复一次反对易，得到
\begin{equation*}
 {
 \Tr(\gamma^\mu\gamma^\nu\gamma^\rho\gamma^\sigma)
 =4(\eta^{\mu\nu}\eta^{\rho\sigma}
 -\eta^{\mu\rho}\eta^{\nu\sigma}
 +\eta^{\mu\sigma}\eta^{\nu\rho}).}
 \end{equation*}
这条式子是后面所有未偏振树级费米子截面的核心；后续未偏振费米子散射的自旋求和将直接使用这一由 Clifford 代数固定的结果。
\end{proof}

\begin{lemma}[偶数 Dirac 迹的递归]
Clifford 关系及迹循环性把 $2n$ 个 gamma 的迹化为 $2n-2$ 个的迹与两两度量缩并之和，其交替符号由移过的 gamma 个数决定，递归式为\eqnrefs{eq:trace-even-recursion-r68}。
\end{lemma}

\begin{proof}
正文\eqnrefs{eq:trace-even-recursion-r68} 已经把 $T_{2n}^{\mu_1\cdots\mu_{2n}}$ 记为长度 $2n$ 的 Dirac 迹。把最左端 $\gamma^{\mu_1}$ 逐次移到右端。第一次交换给

\begin{equation*}
 \gamma^{\mu_1}\gamma^{\mu_2}
 =2\eta^{\mu_1\mu_2}-\gamma^{\mu_2}\gamma^{\mu_1},
\end{equation*}
第二次再对残余的 $\gamma^{\mu_1}\gamma^{\mu_3}$ 使用同一关系。一直进行到第 $2n$ 个矩阵，得到
\begin{align*}
 \gamma^{\mu_1}\cdots\gamma^{\mu_{2n}}
 =&\sum_{j=2}^{2n}(-1)^j 2\eta^{\mu_1\mu_j}
 \gamma^{\mu_2}\cdots\widehat{\gamma^{\mu_j}}\cdots\gamma^{\mu_{2n}}
 \nonumber\\
 &-\gamma^{\mu_2}\cdots\gamma^{\mu_{2n}}\gamma^{\mu_1},
 \end{align*}
其中帽号表示删除该因子；最后一项前的负号来自一共经过 $2n-1$ 次交换。取迹并用迹的循环性，最后一项等于 $-T_{2n}$，因此移到左边得到 $2T_{2n}$。两边除以 $2$：
\begin{equation*}
 {
 T_{2n}^{\mu_1\cdots\mu_{2n}}
 =\sum_{j=2}^{2n}(-1)^j
 \eta^{\mu_1\mu_j}
 \Tr\!\left(
 \gamma^{\mu_2}\cdots\widehat{\gamma^{\mu_j}}\cdots\gamma^{\mu_{2n}}
 \right).}
 \end{equation*}
$n=2$ 时立刻恢复四伽马矩阵迹；$n=3$ 时得到六个 $\gamma$ 矩阵的迹的十五个度规配对。Compton、M\o ller、Bhabha 与圈积分分子中的长矩阵链都可按这一递推逐级缩短。
\end{proof}

\begin{lemma}[单 gamma 的夹心缩并]
在 $d$ 维 Clifford 代数中，$\gamma_\mu\gamma_\alpha\gamma^\mu=(2-d)\gamma_\alpha$。
\end{lemma}

\begin{proof}
第一章的维数正规化已经采用对复维数 $d$ 作解析延拓的约定。若后续积分采用这一调节方法，矩阵代数也必须先保持在同一个 $d$ 维参数下；这里从正文 Clifford 关系直接缩并一个指标，得到正文\eqnrefs{eq:gamma-d-sandwich-r68}。



由 Clifford 代数
\begin{align*}
 \gamma^\mu\gamma^\alpha\gamma_\mu
 &=(2\eta^{\mu\alpha}-\gamma^\alpha\gamma^\mu)\gamma_\mu,\\
 &=2\gamma^\alpha-\gamma^\alpha(\gamma^\mu\gamma_\mu).
\end{align*}
在 $d$ 维中
\begin{equation*}
 \gamma^\mu\gamma_\mu=d\id(),
\end{equation*}
因此
\begin{equation*}
 {
 \gamma^\mu\gamma^\alpha\gamma_\mu=(2-d)\gamma^\alpha.}
 \end{equation*}
四维时就是 $-2\gamma^\alpha$。自能分子的约化直接使用这一恒等式。
\end{proof}

\begin{lemma}[三个 gamma 的夹心缩并]
在同一 $d$ 维约定中，$\gamma_\mu\gamma_\alpha\gamma_\beta\gamma_\gamma\gamma^\mu=-2\gamma_\gamma\gamma_\beta\gamma_\alpha+(4-d)\gamma_\alpha\gamma_\beta\gamma_\gamma$。
\end{lemma}

\begin{proof}
M\o ller、Bhabha、Compton 与圈分子中经常出现把一串 $\gamma$ 矩阵夹在 $\gamma^\mu$ 与 $\gamma_\mu$ 之间的结构。它们都可逐次用

\begin{equation*}
 \gamma^\mu\gamma^\nu=2\eta^{\mu\nu}-\gamma^\nu\gamma^\mu
 \end{equation*}
把左端的 $\gamma^\mu$ 向右移动得到。四维中首先有
\begin{align*}
 \gamma^\mu\slashed a\gamma_\mu
 &=a_\alpha\gamma^\mu\gamma^\alpha\gamma_\mu,\\
 &=a_\alpha(2\gamma^\alpha-\gamma^\alpha\gamma^\mu\gamma_\mu),\\
 &=(2-4)\slashed a=-2\slashed a.
 \end{align*}
对两个斜线记号，
\begin{align*}
 \gamma^\mu\slashed a\slashed b\gamma_\mu
 &=\gamma^\mu\slashed a(2b_\mu-\gamma_\mu\slashed b),\\
 &=2\slashed b\slashed a-
 (\gamma^\mu\slashed a\gamma_\mu)\slashed b,\\
 &=2\slashed b\slashed a+2\slashed a\slashed b,\\
 &=2\{\slashed a,\slashed b\}=4a\!\cdot b\,\id(4).
 \end{align*}
对三个斜线记号，先把最右的 $\slashed c$ 留在外面：
\begin{align*}
 \gamma^\mu\slashed a\slashed b\slashed c\gamma_\mu
 &=2\slashed c\slashed a\slashed b
 -\gamma^\mu\slashed a\slashed b\gamma_\mu\slashed c,\\
 &=2\slashed c\slashed a\slashed b-4(a\!\cdot b)\slashed c.
\end{align*}
再把第一项中的斜线记号重新排序，得到等价形式
\begin{equation*}
 {
 \gamma^\mu\slashed a\slashed b\slashed c\gamma_\mu
 =-2\slashed c\slashed b\slashed a.}
 \end{equation*}
上述四维式是一般 $d$ 维恒等式的特例。由上一引理，在一般维度有
\begin{align*}
\gamma^\mu\slashed a\slashed b\gamma_\mu
&=2\slashed b\slashed a-(2-d)\slashed a\slashed b\\
&=4(a\!\cdot b)I+(d-4)\slashed a\slashed b.
\end{align*}
代入三矩阵夹心，再用
$\slashed c\slashed a\slashed b+\slashed c\slashed b\slashed a=2(a\!\cdot b)\slashed c$，得
\[
\gamma^\mu\slashed a\slashed b\slashed c\gamma_\mu
=-2\slashed c\slashed b\slashed a+(4-d)\slashed a\slashed b\slashed c.
\]
取三个向量分别为坐标基向量即得所述指标形式；四维外腿动量则恢复上面的短矩阵链。
\end{proof}

\begin{lemma}[含手征矩阵的四 gamma 迹]
在四维 Minkowski 度规及本章 $\gamma^5$、$\epsilon^{\mu\nu\rho\sigma}$ 约定下，$\operatorname{tr}(\gamma^5\gamma^\mu\gamma^\nu\gamma^\rho\gamma^\sigma)$ 为完全反对称张量，系数由 $(\mu,\nu,\rho,\sigma)=(0,1,2,3)$ 决定，结果见\eqnrefs{eq:gamma5-four-trace-r68}。此式不能直接作为一般 $d$ 维手征定义。
\end{lemma}

\begin{proof}
轴矢量流和手征结构需要一个与普通四伽马迹不同的对象。正文\eqnrefs{eq:gamma5-def} 已固定 $\gamma^5$ 与取向约定；先证明相应秩四迹完全反对称，再用 $0123$ 分量固定系数，从而得到正文\eqnrefs{eq:gamma5-four-trace-r68}。



固定
\begin{equation*}
 \gamma^5=\ii\gamma^0\gamma^1\gamma^2\gamma^3,
 \qquad \epsilon^{0123}=+1.
\end{equation*}
张量
\begin{equation*}
 T^{\mu\nu\rho\sigma}
 \equiv\Tr(\gamma^5\gamma^\mu\gamma^\nu\gamma^\rho\gamma^\sigma)
\end{equation*}
先证明 $T^{\mu\nu\rho\sigma}$ 的完全反对称性。交换相邻的 $\mu,\nu$ 两个指标时，Clifford 代数给出
\begin{align*}
T^{\nu\mu\rho\sigma}
&=\Tr\!\left(\gamma^5\gamma^\nu\gamma^\mu\gamma^\rho\gamma^\sigma\right),\\
&=-\Tr\!\left(\gamma^5\gamma^\mu\gamma^\nu\gamma^\rho\gamma^\sigma\right)
+2\eta^{\mu\nu}\Tr\!\left(\gamma^5\gamma^\rho\gamma^\sigma\right).
\end{align*}
第二项为零。在当前对角度规的坐标基中，选取不同于 $\rho,\sigma$ 的指标 $\kappa$（不求和）；四维总能作此选择。可逆矩阵 $\gamma^\kappa$ 与 $\gamma^\rho\gamma^\sigma$ 对易，而与 $\gamma^5$ 反对易。因此相似变换给
\begin{align*}
\Tr(\gamma^5\gamma^\rho\gamma^\sigma)
&=\Tr\!\left[\gamma^\kappa(\gamma^5\gamma^\rho\gamma^\sigma)(\gamma^\kappa)^{-1}\right]\\
&=-\Tr(\gamma^5\gamma^\rho\gamma^\sigma)=0.
\end{align*}
于是
$T^{\nu\mu\rho\sigma}=-T^{\mu\nu\rho\sigma}$。对其余任意相邻指标交换重复同一计算，得到 $T$ 对四个 Lorentz 指标完全反对称。四维中任何秩四完全反对称张量都与 $\epsilon^{\mu\nu\rho\sigma}$ 成正比，因此写成
\begin{equation*}
T^{\mu\nu\rho\sigma}=C_5\epsilon^{\mu\nu\rho\sigma}.
\end{equation*}
比例常数 $C_5$ 由 $0123$ 分量确定：
\begin{align*}
 T^{0123}
 &=\Tr(\gamma^5\gamma^0\gamma^1\gamma^2\gamma^3),\\
 &=\Tr\left[
 \ii(\gamma^0\gamma^1\gamma^2\gamma^3)^2
 \right].
\end{align*}
把第二组 $\gamma$ 矩阵逐次穿过第一组并使用
$(\gamma^0)^2=+1$、$(\gamma^i)^2=-1$，得到
\begin{equation*}
 (\gamma^0\gamma^1\gamma^2\gamma^3)^2=-\id(4),
\end{equation*}
所以
\begin{equation*}
 T^{0123}=-4\ii.
\end{equation*}
因此在上述度规与 $\epsilon^{0123}$ 约定下
\begin{equation*}
 {
 \Tr(\gamma^5\gamma^\mu\gamma^\nu\gamma^\rho\gamma^\sigma)
 =-4\ii\epsilon^{\mu\nu\rho\sigma}.}
 \end{equation*}
若改变度规或 $\epsilon^{0123}$ 的符号，右端也必须同步改变；这就是为何约定必须在全书开头冻结。
\end{proof}

\begin{proposition}[Dirac 矩阵的完备基]
四维复 Dirac 矩阵空间由 $I,\gamma^\mu,\sigma^{\mu\nu},\gamma^\mu\gamma^5,\gamma^5$ 的 $1+4+6+4+1=16$ 个独立元素张成。按迹配对取对偶基后，任意四乘四矩阵都有唯一展开\eqnrefs{eq:dirac-matrix-completeness-r10}。
\end{proposition}

\begin{proof}
任意复 $4\times4$ 矩阵形成 $16$ 维线性空间。由 Clifford 代数构造

\begin{equation*}
 \mathcal B=
 \{\id(),\gamma^\mu,\sigma^{\mu\nu},\gamma^\mu\gamma^5,\gamma^5\},
 \end{equation*}
其独立分量数为
$1+4+6+4+1=16$。迹恒等式说明不同 Lorentz 类型之间在迹内积
\begin{equation*}
 (A,B)\equiv\frac14\Tr(AB)
\end{equation*}
下彼此正交；例如 $\Tr\gamma^\mu=0$、$\Tr(\gamma^\mu\gamma^5)=0$、$\Tr\sigma^{\mu\nu}=0$。因此任意矩阵 $M$ 都可唯一展开为
\begin{equation*}
 M=\frac14\sum_A \Gamma_A\Tr(\Gamma^A M),
 \end{equation*}
其中 $A$ 遍历上述 16 个基元素，$\Gamma^A$ 表示按迹度规取对偶基。把 $M$ 取成只有一个元素非零的矩阵单位 $(E_{lk})_{ij}=\delta_{il}\delta_{jk}$，有
\begin{align*}
 \delta_{il}\delta_{jk}
 &=\frac14\sum_A(\Gamma_A)_{ij}
 \Tr(\Gamma^A E_{lk}),\\
 &=\frac14\sum_A(\Gamma_A)_{ij}(\Gamma^A)_{kl}.
\end{align*}
所以得到矩阵指标层面的完备恒等式
\begin{equation*}
 {
 \frac14\sum_A
 (\Gamma_A)_{ij}(\Gamma^A)_{kl}
 =\delta_{il}\delta_{kj}.}
 \end{equation*}
Fierz 重排来自完备矩阵基上的分量展开；它把
\eqnrefs{eq:dirac-matrix-completeness-r10}
插到两个旋量双线性型之间并重新配对旋量指标。若需要某一个具体 Fierz 关系，将从这条母完备关系逐项投影，而不直接引用公式表。
\end{proof}

\begin{proposition}[Weyl 矩阵的行列式]
采用本章 $(+---)$ 度规与 $\sigma$ 矩阵约定，$\det(p_\mu\sigma^\mu)=p_\mu p^\mu$。因此非零二分量旋量满足无质量 Weyl 方程时，其动量位于光锥。
\end{proposition}

\begin{proof}
由正文\eqnrefs{eq:sigma-barsigma-foundation-r10}，四动量被写成 $p\cdot\sigma$ 与 $p\cdot\bar\sigma$ 两个厄米矩阵。这里求它们的两个本征值并相乘，证明行列式正是 Lorentz 不变量 $p^2$，从而得到正文\eqnrefs{eq:weyl-determinant-mass-shell-dp21} 的质量壳条件。



写
\begin{equation*}
 p_\mu\sigma^\mu=p^0I-\mathbf p\cdot\boldsymbol\sigma.
\end{equation*}
对任意 $2\times2$ 矩阵 $M=aI+\mathbf b\cdot\boldsymbol\sigma$，Pauli 代数给
\begin{equation*}
 (\mathbf b\cdot\boldsymbol\sigma)^2=\mathbf b^2I.
\end{equation*}
因此 $\mathbf b\cdot\boldsymbol\sigma$ 的两个本征值是 $\pm|\mathbf b|$，而 $M$ 的本征值是 $a\pm|\mathbf b|$。行列式是本征值乘积：
\begin{equation*}
 \det M=(a+|\mathbf b|)(a-|\mathbf b|)=a^2-\mathbf b^2.
\end{equation*}
取 $a=p^0$、$\mathbf b=-\mathbf p$，得到
\begin{equation*}
 \det(p_\mu\sigma^\mu)
 =(p^0)^2-\mathbf p^2=p_\mu p^\mu.
\end{equation*}
对 $p_\mu\bar\sigma^\mu=p^0I+\mathbf p\cdot\boldsymbol\sigma$ 只把 $\mathbf b$ 改号，行列式不变。

质量壳 $p^2=m^2c^2$ 因而等价于
\begin{equation*}
 \det(p\cdot\sigma)=m^2c^2.
\end{equation*}
这解释了为何有质量 Weyl 方程中 $p\cdot\sigma$ 与 $p\cdot\bar\sigma$ 的乘积恢复二阶质量壳算符 $p^2-m^2c^2$；对电子而言，一阶 Dirac 方程仍是基础动力学方程。
\end{proof}

\begin{proposition}[标量通道的 Fierz 重排]
使用上述 Dirac 基及其迹对偶基，交换两个旋量指标所得恒等式为\eqnrefs{eq:fierz-scalar-r68}。对 Grassmann 奇场重排场算符还需保留交换次序产生的负号。
\end{proposition}

\begin{proof}
正文\eqnrefs{eq:dirac-matrix-completeness-r10} 是 $4\times4$ 矩阵空间的恒等分解。把它插入四个外线旋量的指标之间并改变配对顺序，就能把 $(1,2)(3,4)$ 的标量双线性改写为 $(1,4)(3,2)$ 的五类 Dirac 双线性，得到正文\eqnrefs{eq:fierz-scalar-r68}。



从矩阵完备关系
\begin{equation*}
\delta_{il}\delta_{kj}
=\frac14\sum_A(\Gamma_A)_{ij}(\Gamma^A)_{kl}
\end{equation*}
出发。对四个普通复数值 Dirac 旋量 $u_1,u_2,u_3,u_4$，标量双线性乘积写成显式矩阵指标
\begin{equation*}
(\bar u_1u_2)(\bar u_3u_4)
=(\bar u_1)_i(u_2)_j(\bar u_3)_k(u_4)_l\,
\delta_{ij}\delta_{kl}.
\end{equation*}
配对从 $(1,2)(3,4)$ 交换为 $(1,4)(3,2)$ 时， $\delta_{ij}\delta_{kl}$ 写成 $i$ 与 $l$、$k$ 与 $j$ 相连的基展开。由完备式交换哑指标得
\begin{equation*}
\delta_{ij}\delta_{kl}
=\frac14\sum_A(\Gamma_A)_{il}(\Gamma^A)_{kj}.
\end{equation*}
代入后
\begin{align*}
(\bar u_1u_2)(\bar u_3u_4)
&=\frac14\sum_A
[(\bar u_1)_i(\Gamma_A)_{il}(u_4)_l]
[(\bar u_3)_k(\Gamma^A)_{kj}(u_2)_j]\\
&=\frac14\sum_A
(\bar u_1\Gamma_Au_4)(\bar u_3\Gamma^Au_2).
\end{align*}
现在把 16 维基分成标量、矢量、张量、轴矢量、赝标量五类。按
\begin{equation*}
\sigma^{\mu\nu}=\frac{i}{2}[\gamma^\mu,\gamma^\nu]
\end{equation*}
的归一化，对偶基在张量扇区带因子 $1/2$；轴矢量扇区的迹度规带负号。于是
\begin{align*}
(\bar u_1u_2)(\bar u_3u_4)
=\frac14\Big[&
(\bar u_1u_4)(\bar u_3u_2)
+(\bar u_1\gamma^\mu u_4)(\bar u_3\gamma_\mu u_2)\\
&+\frac12(\bar u_1\sigma^{\mu\nu}u_4)
(\bar u_3\sigma_{\mu\nu}u_2)\\
&-(\bar u_1\gamma^\mu\gamma^5u_4)
(\bar u_3\gamma_\mu\gamma^5u_2)\\
&+(\bar u_1\gamma^5u_4)(\bar u_3\gamma^5u_2)
\Big].
\end{align*}
若四个对象是 Grassmann 量子场，重排 $\psi_2$ 与 $\bar\psi_3$ 到新配对顺序还会产生反对易负号。Fierz 系数由矩阵完备性决定，统计符号则由场算符次序决定。
\end{proof}

\begin{lemma}[两个 Lorentz 自旋生成元的迹]
固定 $\sigma^{\mu\nu}=\frac{\ii}{2}[\gamma^\mu,\gamma^\nu]$，则 $\operatorname{tr}(\sigma^{\mu\nu}\sigma^{\rho\sigma})=4(\eta^{\mu\rho}\eta^{\nu\sigma}-\eta^{\mu\sigma}\eta^{\nu\rho})$。
\end{lemma}

\begin{proof}
正文\eqnrefs{eq:sigma-definition} 已给出 $\sigma^{\mu\nu}$ 与两个 $\gamma$ 矩阵对易子的关系。将两个这样的对象相乘并展开成四伽马迹后，所有张量结构都能由已知迹公式缩并，最终得到正文\eqnrefs{eq:sigma-sigma-trace-r68}。



定义
\begin{equation*}
\sigma^{\mu\nu}=\frac{i}{2}(\gamma^\mu\gamma^\nu-\gamma^\nu\gamma^\mu).
\end{equation*}
所以
\begin{align*}
\Tr(\sigma^{\mu\nu}\sigma^{\rho\sigma})
=-\frac14\Tr\Big[&
\gamma^\mu\gamma^\nu\gamma^\rho\gamma^\sigma
-\gamma^\mu\gamma^\nu\gamma^\sigma\gamma^\rho\\
&-\gamma^\nu\gamma^\mu\gamma^\rho\gamma^\sigma
+\gamma^\nu\gamma^\mu\gamma^\sigma\gamma^\rho
\Big].
\end{align*}
对四项分别代入
\(
\Tr(\gamma^a\gamma^b\gamma^c\gamma^d)
=4(\eta^{ab}\eta^{cd}-\eta^{ac}\eta^{bd}+\eta^{ad}\eta^{bc})
\)。四个迹依次为
\begin{align*}
T_1&=4(\eta^{\mu\nu}\eta^{\rho\sigma}
-\eta^{\mu\rho}\eta^{\nu\sigma}
+\eta^{\mu\sigma}\eta^{\nu\rho}),\\
T_2&=4(\eta^{\mu\nu}\eta^{\rho\sigma}
-\eta^{\mu\sigma}\eta^{\nu\rho}
+\eta^{\mu\rho}\eta^{\nu\sigma}),\\
T_3&=4(\eta^{\mu\nu}\eta^{\rho\sigma}
-\eta^{\nu\rho}\eta^{\mu\sigma}
+\eta^{\nu\sigma}\eta^{\mu\rho}),\\
T_4&=4(\eta^{\mu\nu}\eta^{\rho\sigma}
-\eta^{\nu\sigma}\eta^{\mu\rho}
+\eta^{\nu\rho}\eta^{\mu\sigma}).
\end{align*}
原组合是 $-(T_1-T_2-T_3+T_4)/4$。代入后，四个
$\eta^{\mu\nu}\eta^{\rho\sigma}$ 项的系数为 $1-1-1+1=0$；其余两种张量结构分别收集为
\begin{align*}
-\frac14(T_1-T_2-T_3+T_4)
&=4\eta^{\mu\rho}\eta^{\nu\sigma}
-4\eta^{\mu\sigma}\eta^{\nu\rho}.
\end{align*}
因此
\begin{equation*}
{
\Tr(\sigma^{\mu\nu}\sigma^{\rho\sigma})
=4(\eta^{\mu\rho}\eta^{\nu\sigma}
-\eta^{\mu\sigma}\eta^{\nu\rho}).}
\end{equation*}
进一步令 $\rho=\mu$、$\sigma=\nu$ 并求和，得到
\begin{align*}
\Tr(\sigma^{\mu\nu}\sigma_{\mu\nu})
&=4(\delta^\mu_\mu\delta^\nu_\nu-\delta^\mu_\nu\delta^\nu_\mu)\\
&=4(16-4)=48.
\end{align*}
由于 $\sigma^{\mu\nu}\sigma_{\mu\nu}$ 是 Lorentz 标量并在不可约 Dirac 矩阵空间中与恒等式成正比，而 $\Tr I_4=4$，因此
\begin{equation*}
\sigma^{\mu\nu}\sigma_{\mu\nu}=12I_4.
\end{equation*}
这两个式子是磁矩、Pauli 形状因子和张量算符中最常用的 Clifford 收缩。
\end{proof}

\begin{lemma}[自旋生成元的 Hodge 对偶]
在四维及本章 Levi--Civita 约定下，$\sigma^{\mu\nu}\gamma^5$ 等于 $\sigma_{\rho\sigma}$ 的 Hodge 对偶乘固定虚系数，具体为\eqnrefs{eq:sigma-gamma5-dual-r68}。系数与指标升降必须同时使用同一度规约定。
\end{lemma}

\begin{proof}
在四维和当前约定 $\epsilon^{0123}=+1$ 下，考虑两个秩二反对称张量

\begin{equation*}
\sigma^{\mu\nu}\gamma^5,
\qquad
\epsilon^{\mu\nu\rho\sigma}\sigma_{\rho\sigma}.
\end{equation*}
同一变换类型本身不足以证明成比例，因为复二形式分解为两个手征部分。这里直接由 Clifford 乘法核对六个独立分量。先写待验证关系
\begin{equation*}
\sigma^{\mu\nu}\gamma^5
=C_D\epsilon^{\mu\nu\rho\sigma}\sigma_{\rho\sigma}.
\end{equation*}
取 $(\mu,\nu)=(0,1)$ 来求 $C_D$。由于
$\gamma^1\gamma^0=-\gamma^0\gamma^1$，有
\begin{align*}
\sigma^{01}
&=\frac{i}{2}(\gamma^0\gamma^1-\gamma^1\gamma^0),\\
&=i\gamma^0\gamma^1.
\end{align*}
所以
\begin{align*}
\sigma^{01}\gamma^5
&=(i\gamma^0\gamma^1)
(i\gamma^0\gamma^1\gamma^2\gamma^3),\\
&=-\gamma^0\gamma^1\gamma^0\gamma^1\gamma^2\gamma^3,\\
&=+(\gamma^0)^2(\gamma^1)^2\gamma^2\gamma^3,\\
&=-\gamma^2\gamma^3.
\end{align*}
另一方面，$\epsilon^{01\rho\sigma}$ 只有 $(\rho,\sigma)=(2,3)$ 和 $(3,2)$ 两项：
\begin{align*}
\epsilon^{01\rho\sigma}\sigma_{\rho\sigma}
&=\epsilon^{0123}\sigma_{23}+\epsilon^{0132}\sigma_{32},\\
&=\sigma_{23}+(-1)(-\sigma_{23}),\\
&=2\sigma_{23}.
\end{align*}
对空间指标降指标不改变 $\sigma_{23}$ 的符号，而
\begin{align*}
\sigma_{23}=\sigma^{23}
&=\frac{i}{2}(\gamma^2\gamma^3-\gamma^3\gamma^2),\\
&=i\gamma^2\gamma^3.
\end{align*}
因此右边等于 $2i\gamma^2\gamma^3$。比较
$-\gamma^2\gamma^3=C_D(2i\gamma^2\gamma^3)$ 得到
$C_D=i/2$。其余独立分量逐次反对易约化为
\begin{align*}
\sigma^{02}\gamma^5&=+\gamma^1\gamma^3,&
\sigma^{03}\gamma^5&=-\gamma^1\gamma^2,\\
\sigma^{12}\gamma^5&=+\gamma^0\gamma^3,&
\sigma^{13}\gamma^5&=-\gamma^0\gamma^2,\\
\sigma^{23}\gamma^5&=+\gamma^0\gamma^1.
\end{align*}
对每个右端，$\frac{i}{2}\epsilon^{\mu\nu\rho\sigma}\sigma_{\rho\sigma}$ 给出相同矩阵；这里 $\sigma_{0i}=-\sigma^{0i}$、$\sigma_{ij}=\sigma^{ij}$。反对称性再覆盖反序分量及零对角分量，于是
\begin{equation*}
{
\sigma^{\mu\nu}\gamma^5
=\frac{i}{2}\epsilon^{\mu\nu\rho\sigma}\sigma_{\rho\sigma}.}
\end{equation*}
若翻转度规或 $\epsilon^{0123}$ 约定，右端符号必须同步改变。因此该对偶恒等式的符号由度规与 $\epsilon^{0123}$ 的取向约定共同固定。
\end{proof}

\begin{theorem}[Poincaré Lie 代数]
对平直时空的仿射等距群取单位元处微分，平移生成元彼此交换，Lorentz 生成元按二阶反对称张量的括号闭合，并使平移生成元按向量表示变换；完整关系为\eqnrefs{eq:poincare-mm}及其相邻两式。
\end{theorem}

\begin{proof}
正文已经把 Poincar\'e 群写成仿射矩阵。这里仅完成从矩阵乘法到三组结构常数的代数，并用表示微分同态把抽象基送到厄米量子生成元。

取
\begin{equation*}
 \mathsf x(X,u)=
 \begin{pmatrix}X&u\\0&0\end{pmatrix},
 \qquad
 \mathsf x(Y,v)=
 \begin{pmatrix}Y&v\\0&0\end{pmatrix}.
\end{equation*}
直接相乘得到
\begin{align*}
 \mathsf x(X,u)\mathsf x(Y,v)
 &=\begin{pmatrix}XY&Xv\\0&0\end{pmatrix},\\
 \mathsf x(Y,v)\mathsf x(X,u)
 &=\begin{pmatrix}YX&Yu\\0&0\end{pmatrix},
\end{align*}
故
\begin{equation*}
 [\mathsf x(X,u),\mathsf x(Y,v)]
 =\mathsf x([X,Y],Xv-Yu),
\end{equation*}
即正文\eqnrefs{eq:poincare-semidir-lie-bracket-dp75}。

现在取 Lorentz 基 $m^{\mu\nu}=(\ell^{\mu\nu},0)$ 与平移基 $\tau^\rho=(0,e^\rho)$。两个纯平移元素的矩阵乘积在 Lie 括号中没有非零线性块，所以
\begin{equation*}
 [\tau^\mu,\tau^\nu]=0.
\end{equation*}
Lorentz 基对平移四矢量的作用由正文\eqnrefs{eq:lorentz-matrix-basis} 给出，
\begin{equation*}
 [m^{\mu\nu},\tau^\rho]
 =\eta^{\nu\rho}\tau^\mu-\eta^{\mu\rho}\tau^\nu.
\end{equation*}
而两个 Lorentz 基的矩阵对易子为
\begin{equation*}
[m^{\mu\nu},m^{\rho\sigma}]
=\eta^{\nu\rho}m^{\mu\sigma}
-\eta^{\mu\rho}m^{\nu\sigma}
-\eta^{\nu\sigma}m^{\mu\rho}
+\eta^{\mu\sigma}m^{\nu\rho}.
\end{equation*}
最后使用
\begin{equation*}
 dU(m^{\mu\nu})=-\frac{\ii}{\hbar}M^{\mu\nu},
 \qquad
 dU(\tau^\rho)=-\frac{\ii}{\hbar}P^\rho,
 \qquad
 dU([A,B])=[dU(A),dU(B)],
\end{equation*}
三条抽象 Lie 括号分别变成
\begin{align*}
 [P^\mu,P^\nu]&=0,\\
 [M^{\mu\nu},P^\rho]
 &=\ii\hbar(\eta^{\nu\rho}P^\mu-\eta^{\mu\rho}P^\nu),\\
 [M^{\mu\nu},M^{\rho\sigma}]
 &=\ii\hbar\left(
 \eta^{\nu\rho}M^{\mu\sigma}
 -\eta^{\mu\rho}M^{\nu\sigma}
 -\eta^{\nu\sigma}M^{\mu\rho}
 +\eta^{\mu\sigma}M^{\nu\rho}\right).
\end{align*}
这就是正文\eqnrange{eq:poincare-pp}{eq:poincare-mm}。轨道表达式 $L^{\mu\nu}=x^\mu P^\nu-x^\nu P^\mu$ 只是随后的一种具体表示，不参与一般 Poincar\'e 代数的建立。
\end{proof}

\begin{corollary}[旋转与 boost 的 Lie 括号]
按本章 $J_i,K_i$ 的定义，$[J_i,J_j]=\ii\hbar\epsilon_{ijk}J_k$、$[J_i,K_j]=\ii\hbar\epsilon_{ijk}K_k$、$[K_i,K_j]=-\ii\hbar\epsilon_{ijk}J_k$。
\end{corollary}

\begin{proof}
正文\eqnrefs{eq:poincare-mm} 已给出 $M^{\mu\nu}$ 的协变 Lorentz 代数。把六个独立分量按 $J_i=\epsilon_{ijk}M^{jk}/2$ 与 $K_i=M^{0i}$ 重新组织，并逐项代入协变指标关系，就得到正文\eqnrefs{eq:JK-algebra} 的三组 $J$--$K$ 对易子。


由定义
\begin{equation*}
 J_i=\frac12\epsilon_{ijk}M^{jk},
 \qquad K_i=M^{0i},
\end{equation*}
且反解为 $M^{ij}=\epsilon_{ijk}J_k$。

先取 $(\mu,\nu,\rho,\sigma)=(0,i,0,j)$：
\begin{align*}
 [K_i,K_j]
 &=[M^{0i},M^{0j}],\\
 &=\ii\hbar\left(
 \eta^{i0}M^{0j}-\eta^{00}M^{ij}
 -\eta^{ij}M^{00}+\eta^{0j}M^{i0}\right),\\
 &=-\ii\hbar M^{ij}
 =-\ii\hbar\epsilon_{ijk}J_k.
\end{align*}
这里使用 $\eta^{00}=1$、$\eta^{i0}=0$ 与 $M^{00}=0$。

转动与推动的对易子为
\begin{align*}
 [J_i,K_j]
 &=\frac12\epsilon_{iab}[M^{ab},M^{0j}],\\
 &=\frac{\ii\hbar}{2}\epsilon_{iab}
 (\delta_{bj}M^{a0}-\delta_{aj}M^{b0}),\\
 &=\ii\hbar\epsilon_{ijk}K_k.
\end{align*}

最后对纯空间指标，
\begin{equation*}
 [M^{ab},M^{cd}]
 =\ii\hbar(-\delta_{bc}M^{ad}+\delta_{ac}M^{bd}
 +\delta_{bd}M^{ac}-\delta_{ad}M^{bc}).
\end{equation*}
代入两个 $J$ 并使用
\begin{equation*}
 \epsilon_{iab}\epsilon_{jcb}
 =\delta_{ij}\delta_{ac}-\delta_{ic}\delta_{aj}
\end{equation*}
及其指标置换，可得
\begin{align*}
 [J_i,J_j]
 &=\frac14\epsilon_{iab}\epsilon_{jcd}[M^{ab},M^{cd}],\\
 &=\ii\hbar M^{ij}
 =\ii\hbar\epsilon_{ijk}J_k.
\end{align*}
三组分量关系因此完全由同一个协变 Lorentz 代数得到。
\end{proof}

\begin{proposition}[Pauli--Lubanski 向量的 Lorentz 变换]
由本章定义 $W^\mu$ 构造的 Pauli--Lubanski 算符在 Lorentz 变换下按向量变换，其交换关系为\eqnrefs{eq:m-w}，并满足 $W\cdot P=0$。
\end{proposition}

\begin{proof}
正文把 $W^\mu$ 称为赝四矢量，因此必须验证其四个分量确实按 Lorentz 四矢量混合。对连通保向 Lorentz 变换，$\det\Lambda=1$，完全反对称张量满足
\begin{equation*}
 \epsilon^{\mu\nu\rho\sigma}
 \Lambda_\nu{}^{\alpha}\Lambda_\rho{}^{\beta}\Lambda_\sigma{}^{\gamma}
 =\Lambda^\mu{}_{\delta}\epsilon^{\delta\alpha\beta\gamma}.
\end{equation*}
四动量与 Lorentz 生成元的有限伴随变换分别为
\begin{equation*}
 U^{-1}P_\nu U=\Lambda_\nu{}^{\alpha}P_\alpha,
 \qquad
 U^{-1}M_{\rho\sigma}U
 =\Lambda_\rho{}^{\beta}\Lambda_\sigma{}^{\gamma}M_{\beta\gamma}.
\end{equation*}
于是
\begin{align*}
 U^{-1}W^\mu U
 &=-\frac12\epsilon^{\mu\nu\rho\sigma}
 \bigl(U^{-1}P_\nu U\bigr)
 \bigl(U^{-1}M_{\rho\sigma}U\bigr),\\
 &=-\frac12\epsilon^{\mu\nu\rho\sigma}
 \Lambda_\nu{}^{\alpha}\Lambda_\rho{}^{\beta}\Lambda_\sigma{}^{\gamma}
 P_\alpha M_{\beta\gamma},\\
 &=\Lambda^\mu{}_{\delta}W^\delta.
\end{align*}
再令 $\Lambda^\mu{}_{\nu}=\delta^\mu{}_{\nu}+\omega^\mu{}_{\nu}+\order(\omega^2)$，并展开
\begin{equation*}
 U=\id()-\frac{\ii}{2\hbar}\omega_{\alpha\beta}M^{\alpha\beta}
 +\order(\omega^2),
\end{equation*}
比较任意反对称参数 $\omega_{\alpha\beta}$ 的一阶系数，得到
\begin{equation*}
 [M^{\alpha\beta},W^\mu]
 =\ii\hbar\left(\eta^{\beta\mu}W^\alpha-
 \eta^{\alpha\mu}W^\beta\right),
\end{equation*}
即正文\eqnrefs{eq:m-w}。若把宇称也包括进不连通的 Lorentz 群，\eqnrefs{eq:pauli-lubanski} 中的 $\epsilon^{\mu\nu\rho\sigma}$ 使 $W^\mu$ 具有赝矢量的额外宇称符号。
\end{proof}

\begin{proposition}[Pauli--Lubanski 算符的闭合关系]
同一 Poincaré 表示中，$[W^\mu,P^\nu]=0$，两个 $W$ 的交换子由 Levi--Civita 张量、$W$ 与 $P$ 的乘积给出，符号按\eqnrefs{eq:w-w-dp34}。
\end{proposition}

\begin{proof}
直接从 $W^\mu$ 的定义出发，并使用 $[P_\alpha,W^\nu]=0$：
\begin{align*}
[W^\mu,W^\nu]
&=-\frac12\epsilon^{\mu\alpha\rho\sigma}
[P_\alpha M_{\rho\sigma},W^\nu]\\
&=-\frac12\epsilon^{\mu\alpha\rho\sigma}
P_\alpha[M_{\rho\sigma},W^\nu].
\end{align*}
由\eqnrefs{eq:m-w} 降低生成元指标，
\begin{equation*}
[M_{\rho\sigma},W^\nu]
=\ii\hbar(\delta_\sigma^\nu W_\rho-\delta_\rho^\nu W_\sigma).
\end{equation*}
代回并合并两个反对称项：
\begin{align*}
[W^\mu,W^\nu]
&=\ii\hbar\,\epsilon^{\mu\alpha\nu\rho}P_\alpha W_\rho\\
&=\ii\hbar\,\epsilon^{\mu\nu\rho\sigma}W_\rho P_\sigma,
\end{align*}
其中最后一步只重排 Levi--Civita 指标并使用 $[W_\rho,P_\sigma]=0$。这就是正文\eqnrefs{eq:w-w-dp34}，整个推导没有选择静止系或特定质量轨道。
\end{proof}

\begin{theorem}[Poincaré 二次 Casimir]
$P^2$ 与 $W^2$ 均与所有 Poincaré 生成元交换。有质量不可约幺正表示的静止系中，$W^2=-m^2c^2\hbar^2s(s+1)$，并据此确定自旋 $s$。
\end{theorem}

\begin{proof}
Casimir 算符的定义不是“在某个参考系里数值不变”，而是它与整个 Lie 代数的所有生成元对易。先证明
\[
C_1\equiv P_\mu P^\mu
\]
满足这一要求。平移生成元彼此对易，所以
\[
[P^\lambda,C_1]=0.
\]
再用 Poincar\'e 代数
\[
[M^{\alpha\beta},P^\mu]
=\ii\hbar
\left(
\eta^{\beta\mu}P^\alpha
-\eta^{\alpha\mu}P^\beta
\right)
\]
计算
\begin{align*}
[M^{\alpha\beta},P_\mu P^\mu]
&=[M^{\alpha\beta},P_\mu]P^\mu
+P_\mu[M^{\alpha\beta},P^\mu]\\
&=\ii\hbar
\left(
P^\alpha P^\beta-P^\beta P^\alpha
\right)
+\ii\hbar
\left(
P^\beta P^\alpha-P^\alpha P^\beta
\right)\\
&=0.
\end{align*}
因此 \(P^2\) 与全部平移、Lorentz 生成元对易，是第一个 Casimir。

第二个 Casimir 从 Pauli--Lubanski 四矢量
\[
W^\mu
\equiv
-\frac12\epsilon^{\mu\nu\rho\sigma}
P_\nu M_{\rho\sigma}
\]
构造。先看它与平移的对易子：
\begin{align*}
[P^\lambda,W^\mu]
&=-\frac12\epsilon^{\mu\nu\rho\sigma}
P_\nu[P^\lambda,M_{\rho\sigma}]\\
&=+\frac{\ii\hbar}{2}
\epsilon^{\mu\nu\rho\sigma}
P_\nu
\left(
\delta^\lambda_{\ \sigma}P_\rho
-\delta^\lambda_{\ \rho}P_\sigma
\right).
\end{align*}
每一项都含对称乘积 \(P_\nu P_\rho\) 与完全反对称的
\(\epsilon^{\mu\nu\rho\sigma}\) 收缩，因此
\[
[P^\lambda,W^\mu]=0.
\]
同样，从 \([M,M]\) 与 \([M,P]\) 两部分逐项展开可得
\[
[M^{\alpha\beta},W^\mu]
=\ii\hbar
\left(
\eta^{\beta\mu}W^\alpha
-\eta^{\alpha\mu}W^\beta
\right),
\]
即 \(W^\mu\) 确实按 Lorentz 四矢量变换。于是
\begin{align*}
[M^{\alpha\beta},W_\mu W^\mu]
&=[M^{\alpha\beta},W_\mu]W^\mu
+W_\mu[M^{\alpha\beta},W^\mu]\\
&=0,
\end{align*}
而 \([P^\lambda,W^2]=0\) 已由 \([P,W]=0\) 立即得到。因此
\[
C_2\equiv W_\mu W^\mu
\]
也是 Poincar\'e 代数的 Casimir。

还可直接验证
\[
W^\mu P_\mu
=
-\frac12
\epsilon^{\mu\nu\rho\sigma}
P_\mu P_\nu M_{\rho\sigma}
=0,
\]
因为 \(P_\mu P_\nu\) 对 \(\mu,\nu\) 对称。对有质量不可约表示 \(m>0\)，选择静止系
\[
P^\mu=(mc,\mathbf0).
\]
正文\eqnrefs{eq:W-rest-r13} 给出
\[
W^0=0,
\qquad
\mathbf W=mc\,\mathbf J_{\rm int},
\]
其中 \(\mathbf J_{\rm int}\) 是静止小群 \(SO(3)\) 的内部角动量生成元。于是
\[
W^2
=(W^0)^2-\mathbf W^2
=-m^2c^2\,\mathbf J_{\rm int}^2.
\]
不可约自旋 \(s\) 表示满足
\[
\mathbf J_{\rm int}^2
=\hbar^2s(s+1),
\]
所以
\[
\boxed{
P^2=m^2c^2,
\qquad
W^2=-m^2c^2\hbar^2s(s+1).
}
\]
这就是正文\eqnrefs{eq:poincare-casimirs}。质量与自旋之所以能标记 Poincar\'e 不可约表示，不是由静止系经验“读出来”的，而是因为对应算符先由代数证明为 Casimir，再在静止系评价其本征值。
\end{proof}

\begin{lemma}[Pauli 乘积与 Levi--Civita 缩并]
Pauli 矩阵满足 $\sigma_i\sigma_j=\delta_{ij}I+\ii\epsilon_{ijk}\sigma_k$。对应三维完全反对称符号满足 $\epsilon_{ijk}\epsilon_{imn}=\delta_{jm}\delta_{kn}-\delta_{jn}\delta_{km}$，其他缩并由重标指标得到。
\end{lemma}

\begin{proof}
Pauli 乘法\eqnrefs{eq:pauli-product-proof-r10} 的反对称部分由 $\epsilon_{ijk}$ 编码；两个这样的结构相乘时，需要把 Levi--Civita 的双重缩并化成 Kronecker $\delta$。这里从三维叉积的分量定义出发，补出正文\eqnrefs{eq:epsilon-contractions-r10} 的指标恒等式。



取标准正交基 $\mathbf e_i$，由
\begin{equation*}
 \epsilon_{ijk}=\mathbf e_i\cdot(\mathbf e_j\times\mathbf e_k)
\end{equation*}
以及完备关系 $\sum_i(\mathbf e_i\cdot\mathbf a)(\mathbf e_i\cdot\mathbf b)=\mathbf a\cdot\mathbf b$，有
\begin{align*}
 \epsilon_{ijk}\epsilon_{ilm}
 &= (\mathbf e_j\times\mathbf e_k)\cdot(\mathbf e_l\times\mathbf e_m)\\
 &= (\mathbf e_j\cdot\mathbf e_l)(\mathbf e_k\cdot\mathbf e_m)
 -(\mathbf e_j\cdot\mathbf e_m)(\mathbf e_k\cdot\mathbf e_l)\\
 &=\delta_{jl}\delta_{km}-\delta_{jm}\delta_{kl}.
\end{align*}
再令 $l=j$ 并对 $j$ 求和，
\begin{align*}
 \epsilon_{ijk}\epsilon_{ijm}
 &=\delta_{jj}\delta_{km}-\delta_{jm}\delta_{kj}\\
 &=3\delta_{km}-\delta_{km}=2\delta_{km},
\end{align*}
得到\eqnrefs{eq:epsilon-contractions-r10}。
\end{proof}

\begin{lemma}[两种 Weyl 矩阵的 Clifford 恒等式]
在 $(+---)$ 度规中，$\sigma^\mu=(I,\boldsymbol\sigma)$、$\bar\sigma^\mu=(I,-\boldsymbol\sigma)$ 满足 $\sigma^\mu\bar\sigma^\nu+\sigma^\nu\bar\sigma^\mu=2\eta^{\mu\nu}I$，其对换版本也成立。
\end{lemma}

\begin{proof}
正文\eqnrefs{eq:sigma-barsigma-foundation-r10} 把四矢量嵌入两组 $2\times2$ Weyl 矩阵。为了使左右手一阶算符的乘积恢复 Minkowski 度规，逐分量检查时间--时间、时间--空间和空间--空间三类组合，便得到正文\eqnrefs{eq:sigmabar-proof-r10}。



由\eqnrefs{eq:sigma-barsigma-foundation-r10} 分三类指标检查。若 $\mu=\nu=0$，
\begin{equation*}
 \sigma^0\bar\sigma^0+\sigma^0\bar\sigma^0=2\id(2).
\end{equation*}
若一个指标为时间、另一个为空间，
\begin{equation*}
 \sigma^0\bar\sigma^i+\sigma^i\bar\sigma^0
 =-\sigma_i+\sigma_i=0.
\end{equation*}
若两个指标都为空间，\eqnrefs{eq:pauli-comm-anticomm-r10} 给
\begin{align*}
 \sigma^i\bar\sigma^j+\sigma^j\bar\sigma^i
 &=-(\sigma_i\sigma_j+\sigma_j\sigma_i)\\
 &=-2\delta_{ij}\id(2).
\end{align*}
三种情况分别对应 $\eta^{00}=1$、$\eta^{0i}=0$、$\eta^{ij}=-\delta_{ij}$，故合并为\eqnrefs{eq:sigmabar-proof-r10}。
\end{proof}

\begin{theorem}[SL(2,C) 对保向保时 Lorentz 群的双覆盖]
作用 $X=x^0I_2+\mathbf x\cdot\boldsymbol\sigma_{\rm P}\mapsto AXA^\dagger$，$A\in SL(2,\CC)$，诱导到 $SO^+(1,3)$ 的满同态，核为 $\{I,-I\}$。因而 $SO^+(1,3)\simeq SL(2,\CC)/\{\pm I\}$。沿本章已固定的 $\bar\sigma^\mu=(I,-\boldsymbol\sigma_{\rm P})$，这里也可写 $X=x_\mu\bar\sigma^\mu$。
\end{theorem}

\begin{proof}
先计算

\begin{align*}
 \det X
 &=(x^0+x^3)(x^0-x^3)
 -(x^1-\ii x^2)(x^1+\ii x^2)\\
 &=(x^0)^2-(x^1)^2-(x^2)^2-(x^3)^2\\
 &=x_\mu x^\mu.
\end{align*}
对 $A\in SL(2,\mathbb C)$，$X'=AXA^\dagger$ 仍为厄米，且
\begin{align*}
 \det X'
 &=\det A\,\det X\,\det A^\dagger\\
 &=|\det A|^2\det X\\
 &=\det X.
\end{align*}
因此存在唯一实线性变换 $\Phi(A)$ 使 $X(x')=AX(x)A^\dagger$，并且 $\Phi(A)$ 保持 Minkowski 度规。矩阵乘法还给出
\begin{equation*}
 \Phi(A_2A_1)X
 =A_2A_1X A_1^\dagger A_2^\dagger
 =\Phi(A_2)\Phi(A_1)X,
\end{equation*}
所以 $\Phi$ 是群同态。由于 $SL(2,\mathbb C)$ 连通且 $\Phi(\id(2))=\id(4)$，其像落在 Lorentz 群的 proper orthochronous 连通分支 $SO^+(1,3)$。

严格求核。若 $A\in\ker\Phi$，则对任意厄米 $X$ 都有
\begin{equation*}
 AXA^\dagger=X.
\end{equation*}
先取 $X=\id(2)$，得到 $AA^\dagger=\id(2)$，所以 $A$ 是幺正矩阵，$A^\dagger=A^{-1}$。于是上式改写成
\begin{equation*}
 AXA^{-1}=X
 \qquad\text{对任意厄米 }X.
\end{equation*}
厄米矩阵张成整个 $M_2(\mathbb C)$ 复矩阵空间，因此 $A$ 与所有 $2\times2$ 复矩阵对易，只能是标量矩阵 $A=\lambda\id(2)$。再用 $\det A=\lambda^2=1$，得到
\begin{equation*}
 \ker\Phi=\{+\id(2),-\id(2)\}.
\end{equation*}

最后证明满射，而不是只验证若干特殊变换。把 $\mathfrak{sl}(2,\mathbb C)$ 看成六维实 Lie 代数，其一组实基可取
\begin{equation*}
 -\frac{\ii}{2}\sigma_i,
 \qquad
 +\frac12\sigma_i,
 \qquad i=1,2,3.
\end{equation*}
对转动型无穷小元
\begin{equation*}
 A=\id(2)-\frac{\ii}{2}\boldsymbol\theta\cdot\boldsymbol\sigma_{\rm P}
 +O(\theta^2),
\end{equation*}
有
\begin{align*}
 \delta X
 &=-\frac{\ii}{2}
 [\boldsymbol\theta\cdot\boldsymbol\sigma_{\rm P},
   \mathbf x\cdot\boldsymbol\sigma_{\rm P}]\\
 &=(\boldsymbol\theta\times\mathbf x)\cdot\boldsymbol\sigma_{\rm P},
\end{align*}
即 $\delta x^0=0$、$\delta\mathbf x=\boldsymbol\theta\times\mathbf x$，给出三个独立空间转动方向。对推动型无穷小元
\begin{equation*}
 A=\id(2)+\frac12\boldsymbol\xi\cdot\boldsymbol\sigma_{\rm P}
 +O(\xi^2),
\end{equation*}
则
\begin{align*}
 \delta X
 &=\frac12\{\boldsymbol\xi\cdot\boldsymbol\sigma_{\rm P},X\}\\
 &=(\boldsymbol\xi\cdot\mathbf x)\id(2)
 +x^0\boldsymbol\xi\cdot\boldsymbol\sigma_{\rm P},
\end{align*}
因此 $\delta x^0=\boldsymbol\xi\cdot\mathbf x$、$\delta\mathbf x=x^0\boldsymbol\xi$，给出三个独立 Lorentz 推动方向。故微分
\begin{equation*}
 \dd\Phi_e:\mathfrak{sl}(2,\mathbb C)_{\mathbb R}
 \longrightarrow\mathfrak{so}(1,3)
\end{equation*}
在六个实维度上满秩，因而是 Lie 代数同构。于是 $\Phi(SL(2,\mathbb C))$ 是 $SO^+(1,3)$ 的开子群；而 $SO^+(1,3)$ 连通，开子群的不同陪集若多于一个会把群分成互不相交的非空开集，所以只能有一个陪集。于是 $\Phi$ 满射。

综上，$\Phi$ 是满射群同态，核为 $\{\pm\id(2)\}$，由第一同构定理得到
\begin{equation*}
 SL(2,\mathbb C)/\{\pm\id(2)\}
 \cong SO^+(1,3),
\end{equation*}
即正文\eqnrefs{eq:sl2c-double-cover-dp8}。这里的“双覆盖”因此同时包含“恰有二元核”和“确实覆盖整个 $SO^+(1,3)$”两层结论。
\end{proof}

\begin{proposition}[Wigner 态的 Lorentz 变换]
固定标准动量、把它送到 $p$ 的标准变换以及本章动量态测度归一化。诱导表示中的态在 Lorentz 变换下由小群元 $L^{-1}(\Lambda p)\Lambda L(p)$ 作用于内部指标，并带有\eqnrefs{eq:wigner-state-transform-dp68}中的测度补偿因子。
\end{proposition}

\begin{proof}
对任意保向正时向 Lorentz 变换 $\Lambda$，定义

\begin{equation*}
 {
 W(\Lambda,p)
 \equiv L^{-1}(\Lambda p)\,\Lambda\,L(p).}
 \end{equation*}
直接作用在标准动量上：
\begin{align*}
 W(\Lambda,p)\bar p
 &=L^{-1}(\Lambda p)\Lambda L(p)\bar p,\\
 &=L^{-1}(\Lambda p)\Lambda p,\\
 &=\bar p.
 \end{align*}
因此 $W(\Lambda,p)$ 确实属于稳定 $\bar p$ 的小群
\begin{equation*}
 G_{\bar p}\equiv\{W\in SO^+(1,3)\mid W\bar p=\bar p\}.
 \end{equation*}
再把 \eqnrefs{eq:wigner-state-definition-dp9} 代入态变换：
\begin{align*}
 U(\Lambda)|p,\sigma\rangle
 &=U(\Lambda)U[L(p)]|\bar p,\sigma\rangle,\\
 &=U[L(\Lambda p)]U[W(\Lambda,p)]|\bar p,\sigma\rangle.
\end{align*}
因为 $W$ 不改变 $\bar p$，它只能在固定动量的内部空间中作用。若该内部空间取小群不可约表示 $D$，
\begin{equation*}
 U(W)|\bar p,\sigma\rangle
 =\sum_{\sigma'}D_{\sigma'\sigma}(W)|\bar p,\sigma'\rangle,
\end{equation*}
于是
\begin{equation*}
 {
 U(\Lambda)|p,\sigma\rangle
 =\sum_{\sigma'}
 D_{\sigma'\sigma}[W(\Lambda,p)]
 |\Lambda p,\sigma'\rangle.}
 \end{equation*}
这条式子把四动量的 Lorentz 变换与内部自旋/螺旋度变换严格分离。若改用非协变三维 $\delta$ 归一化，右端还会出现由雅可比行列式产生的 $\sqrt{(\Lambda p)^0/p^0}$ 因子；后续场量子化采用协变归一化，因此不再引入这一额外因子。
\end{proof}

\begin{theorem}[有质量小群]
对于标准类时动量 $k=(mc,\mathbf0)$，$SO^+(1,3)$ 中保持 $k$ 的子群同构于 $SO(3)$，在旋量覆盖群中的原像为 $SU(2)$。
\end{theorem}

\begin{proof}
有质量粒子可选静止标准四动量
\begin{equation*}
 \bar p_m^\mu=(mc,0,0,0).
\end{equation*}
小群元素 $W$ 的定义是 $W\bar p_m=\bar p_m$。逐分量写出
\begin{equation*}
 W^\mu{}_{\nu}\bar p_m^\nu
 =mc\,W^\mu{}_0
 =mc\,\delta^\mu{}_0,
\end{equation*}
所以第一列被完全固定：
\begin{equation*}
 W^0{}_0=1,
 \qquad
 W^i{}_0=0.
\end{equation*}
还需要证明第一行也没有 boost 分量。Lorentz 条件
$W^{\mathsf T}\eta W=\eta$ 的 $(0,i)$ 分量给
\begin{equation*}
 \eta_{\alpha\beta}W^\alpha{}_0W^\beta{}_i=0.
\end{equation*}
代入 $W^\alpha{}_0=\delta^\alpha{}_0$ 后得到
$W^0{}_i=0$。因此 $W$ 必须具有块形式
\begin{equation*}
 W=
 \begin{pmatrix}
 1&\mathbf0^{\mathsf T}\\
 \mathbf0&R
 \end{pmatrix}.
\end{equation*}
再把这个块矩阵代回 Lorentz 条件，空间块满足
\begin{equation*}
 R^{\mathsf T}R=I_3.
\end{equation*}
若只考虑与恒等变换连续相连的 proper orthochronous Lorentz 群，还要求
$\det W=+1$，因此 $\det R=+1$。所以有质量标准动量的连续稳定子正是
$SO(3)$。

量子态的自旋不是用 $SO(3)$ 的普通表示分类，而是用其双覆盖 $SU(2)$ 的幺正不可约表示分类，因为 $2\pi$ 空间转动可以在态矢量上产生负号而不改变物理射线。$SU(2)$ 不可约表示由
$s=0,\tfrac12,1,\ldots$ 标记，维数为 $2s+1$；这正是正文\eqnrefs{eq:wigner-massive-little-group-dp68} 中出现的自旋量子数，也与 Pauli--Lubanski Casimir 的标签一致。
\end{proof}

\begin{theorem}[无质量小群的 Lie 代数]
对于非零未来光状标准动量，其稳定子的 Lie 代数同构于二维 Euclidean 运动代数 $\mathfrak{iso}(2)$。两个平移型生成元彼此交换，绕传播轴的转动旋转这两个生成元，关系为\eqnrefs{eq:wigner-massless-little-group-algebra-dp68}。
\end{theorem}

\begin{proof}
稳定子代数的完备性由一般 Lorentz 代数元的线性稳定条件决定。先解出全部稳定方向，再计算所得三维解空间的 Lie 括号。

写
\begin{equation*}
 X=a_1J_1+a_2J_2+a_3J_3
 +b_1K_1+b_2K_2+b_3K_3.
\end{equation*}
由正文\eqnrefs{eq:poincare-mp}，
\begin{equation*}
 [J_i,P^0]=0,
 \qquad
 [J_i,P^j]=\ii\hbar\epsilon_{ijk}P^k,
\end{equation*}
以及
\begin{equation*}
 [K_i,P^0]=-\ii\hbar P^i,
 \qquad
 [K_i,P^j]=-\ii\hbar\delta_{ij}P^0.
\end{equation*}
在
\begin{equation*}
 \bar k^\mu=(\kappa_\star,0,0,\kappa_\star)
\end{equation*}
上逐分量要求 $[X,P^\mu]_{P=\bar k}=0$。时间分量给
\begin{equation*}
 -\ii\hbar\kappa_\star b_3=0
 \quad\Longrightarrow\quad b_3=0.
\end{equation*}
第一空间分量给
\begin{equation*}
 -\ii\hbar\kappa_\star(a_2+b_1)=0
 \quad\Longrightarrow\quad b_1=-a_2,
\end{equation*}
第二空间分量给
\begin{equation*}
 \ii\hbar\kappa_\star(a_1-b_2)=0
 \quad\Longrightarrow\quad b_2=a_1.
\end{equation*}
第三空间分量在这些条件下自动为零。因此任意稳定子元都唯一写成
\begin{equation*}
 X=a_3J_3+a_1(K_2+J_1)-a_2(K_1-J_2),
\end{equation*}
证明稳定子空间恰为
\begin{equation*}
 \operatorname{span}\{J_3,\mathsf N_1,\mathsf N_2\}.
\end{equation*}

再用正文\eqnrefs{eq:JK-algebra}：
\begin{align*}
 [J_3,\mathsf N_1]
 &=[J_3,K_1]-[J_3,J_2]
 =\ii\hbar(K_2+J_1)
 =\ii\hbar\mathsf N_2,\\
 [J_3,\mathsf N_2]
 &=[J_3,K_2]+[J_3,J_1]
 =-\ii\hbar(K_1-J_2)
 =-\ii\hbar\mathsf N_1,\\
 [\mathsf N_1,\mathsf N_2]
 &=[K_1,K_2]+[K_1,J_1]-[J_2,K_2]-[J_2,J_1]\\
 &=-\ii\hbar J_3+\ii\hbar J_3\\
 &=0.
\end{align*}
这就是正文\eqnrefs{eq:wigner-massless-little-group-algebra-dp68}。由于三维解空间已经由一般稳定条件完整求出，所得 $\mathfrak{iso}(2)$ 不再只是三个候选生成元的子代数，而是零质量标准动量的全部连续稳定子 Lie 代数。
\end{proof}

\begin{theorem}[复化 Lorentz 代数的手征分解]
按本章含 $\hbar$ 的 $J,K$ 约定，令 $A_i=(J_i+\ii K_i)/2$、$B_i=(J_i-\ii K_i)/2$。则 $[A_i,A_j]=\ii\hbar\epsilon_{ijk}A_k$、$[B_i,B_j]=\ii\hbar\epsilon_{ijk}B_k$、$[A_i,B_j]=0$。此分解针对复化代数，实 Lorentz 代数并非两个紧实代数的直和。
\end{theorem}

\begin{proof}
正文\eqnrefs{eq:JK-algebra} 已把六个 Lorentz 生成元写成三维转动 $J_i$ 与推动 $K_i$。将它们组合成 $\mathsf J_i^{(\pm)}=(J_i\pm\ii K_i)/2$ 后，逐项代入三组 $J$--$K$ 对易关系即可证明两组生成元各自闭合且彼此对易，从而得到正文\eqnrefs{eq:lorentz-su2-algebra-dp34}。



由
\begin{equation*}
 \mathsf J_i^{(+)}=\frac12(J_i+\ii K_i),
 \qquad
 \mathsf J_i^{(-)}=\frac12(J_i-\ii K_i)
\end{equation*}
与
\begin{equation*}
 [J_i,J_j]=\ii\hbar\epsilon_{ijk}J_k,
 \qquad
 [J_i,K_j]=\ii\hbar\epsilon_{ijk}K_k,
 \qquad
 [K_i,K_j]=-\ii\hbar\epsilon_{ijk}J_k,
\end{equation*}
先计算第一组：
\begin{align*}
 [\mathsf J_i^{(+)},\mathsf J_j^{(+)}]
 &=\frac14\Bigl(
 [J_i,J_j]+\ii[J_i,K_j]
 +\ii[K_i,J_j]-[K_i,K_j]\Bigr).
\end{align*}
利用
\begin{equation*}
 [K_i,J_j]=-[J_j,K_i]
 =\ii\hbar\epsilon_{ijk}K_k,
\end{equation*}
得到
\begin{align*}
 [\mathsf J_i^{(+)},\mathsf J_j^{(+)}]
 &=\frac14\left(
 2\ii\hbar\epsilon_{ijk}J_k
 -2\hbar\epsilon_{ijk}K_k\right),\\
 &=\ii\hbar\epsilon_{ijk}
 \frac{J_k+\ii K_k}{2},\\
 &=\ii\hbar\epsilon_{ijk}\mathsf J_k^{(+)}.
\end{align*}
第二组从定义 $\mathsf J_i^{(-)}=(J_i-\ii K_i)/2$ 独立计算：
\begin{align*}
 [\mathsf J_i^{(-)},\mathsf J_j^{(-)}]
 &=\frac14\Bigl(
 [J_i,J_j]-\ii[J_i,K_j]
 -\ii[K_i,J_j]-[K_i,K_j]\Bigr),\\
 &=\frac14\Bigl(
 \ii\hbar\epsilon_{ijk}J_k
 +\hbar\epsilon_{ijk}K_k
 +\hbar\epsilon_{ijk}K_k
 +\ii\hbar\epsilon_{ijk}J_k\Bigr),\\
 &=\frac12\Bigl(
 \ii\hbar\epsilon_{ijk}J_k
 +\hbar\epsilon_{ijk}K_k\Bigr),\\
 &=\ii\hbar\epsilon_{ijk}\frac{J_k-\ii K_k}{2},\\
 &=\ii\hbar\epsilon_{ijk}\mathsf J_k^{(-)}.
\end{align*}
交叉对易子为
\begin{align*}
 [\mathsf J_i^{(+)},\mathsf J_j^{(-)}]
 &=\frac14\Bigl(
 [J_i,J_j]-\ii[J_i,K_j]
 +\ii[K_i,J_j]+[K_i,K_j]\Bigr),\\
 &=0.
\end{align*}
因此复化 Lorentz 李代数分解为两份彼此对易的 $\mathfrak{sl}(2,\mathbb C)$；这里只是采用了角动量型的复基，并未把非紧致实 Lorentz 代数写成紧致实代数的直和。
\end{proof}

\begin{proposition}[有质量 Weyl 分量的耦合]
在本章 Weyl 基中，四分量 Dirac 方程等价于左右两个 Weyl 分量的耦合方程\eqnrefs{eq:massive-weyl-coupled-r7}。质量非零时单个手征分量不闭合；质量为零时两方程解耦。
\end{proposition}

\begin{proof}
正文已指出 $p_\mu\sigma^\mu$ 与 $p_\mu\bar\sigma^\mu$ 是两个手征表示之间的一阶 Lorentz 交织映射。这里只需确定允许的标量系数，并由质量壳固定它们的乘积；静止系不参与一般式的建立。

先写最一般线性耦合
\begin{equation*}
 (p_\mu\sigma^\mu)\psi_R=\mu_L\psi_L,
 \qquad
 (p_\mu\bar\sigma^\mu)\psi_L=\mu_R\psi_R.
\end{equation*}
由正文\eqnrefs{eq:sigmabar-identity}，
\begin{align*}
 (p_\mu\bar\sigma^\mu)(p_\nu\sigma^\nu)
 &=\frac12p_\mu p_\nu
 \left(\bar\sigma^\mu\sigma^\nu+\bar\sigma^\nu\sigma^\mu\right)\\
 &=p_\mu p^\mu\,\id(2).
\end{align*}
把第一式左乘 $p_\mu\bar\sigma^\mu$，再代入第二式：
\begin{align*}
 p^2\psi_R
 &=(p\!\cdot\!\bar\sigma)(p\!\cdot\!\sigma)\psi_R\\
 &=\mu_L(p\!\cdot\!\bar\sigma)\psi_L\\
 &=\mu_L\mu_R\psi_R.
\end{align*}
对非零态得到
\begin{equation*}
 p^2=\mu_L\mu_R.
\end{equation*}
同样作用在 $\psi_L$ 上得到相同条件。对质量壳 $p^2=m_{\mathrm e}^2c^2$，
\begin{equation*}
 \mu_L\mu_R=m_{\mathrm e}^2c^2.
\end{equation*}

若自由二次作用量厄米，左右两条方程中的质量系数互为复共轭，
\begin{equation*}
 \mu_R=\mu_L^*.
\end{equation*}
于是 $|\mu_L|=|\mu_R|=m_{\mathrm e} c$。对两个手征场作相反的常数重相位可消去这一个共同复相位，因此可选
\begin{equation*}
 \mu_L=\mu_R=m_{\mathrm e} c>0.
\end{equation*}
这就得到正文\eqnrefs{eq:massive-weyl-coupled-r7}。

作为随后才做的检查，把质量壳写成正文\eqnrefs{eq:rapidity-momentum-r7}。由 Pauli 代数，
\begin{align*}
 p_\mu\sigma^\mu
 &=m_{\mathrm e} c\left(\cosh\zeta\,\id(2)
 -\sinh\zeta\,\hat{\mathbf n}\cdot\boldsymbol\sigma_{\rm P}\right)\\
 &=m_{\mathrm e} c\,\ee^{-\zeta\hat{\mathbf n}\cdot\boldsymbol\sigma_{\rm P}},\\
 p_\mu\bar\sigma^\mu
 &=m_{\mathrm e} c\,\ee^{+\zeta\hat{\mathbf n}\cdot\boldsymbol\sigma_{\rm P}}.
\end{align*}
与正文\eqnrefs{eq:weyl-boost-r7} 的左右手有限推动一致。快度表示在这里提供一般耦合式的一种显式质量壳参数化检验；一般方程本身已经由 Lorentz 交织映射与质量壳条件确定。
\end{proof}

\begin{theorem}[四维复 Clifford 模的最小维数]
四维非退化复 Clifford 代数同构于 $M_4(\CC)$，其非零不可约模维数为四。一个二分量 Weyl 模仅是偶子代数的不可约模，不能单独承载全部四个 gamma。
\end{theorem}

\begin{proof}
设 $A,B$ 是任意两个不同的 Clifford 生成元，因此

\begin{equation*}
 AB=-BA,
 \qquad
 \det A\neq0,
 \qquad
 \det B\neq0.
\end{equation*}
若表示维数为奇数 $n$，则
\begin{align*}
 \det(AB)
 &=\det(-BA)\\
 &=(-1)^n\det(BA)\\
 &=-\det(AB),
\end{align*}
与 $\det(AB)\neq0$ 矛盾。因此任何有限维复 Clifford 表示的维数在这里都不能为奇数；特别地，三维表示已经被排除。

还需排除唯一更小的偶数候选 $n=2$。任意 $2\times2$ 复矩阵可唯一写成
\begin{equation*}
 X=a\id(2)+\mathbf b\cdot\boldsymbol\sigma_{\rm P}.
\end{equation*}
若三个空间 Clifford 生成元存在于二维空间，经相似变换与非零常数重标度后可取为三条 Pauli 方向。与三者都反对易的时间生成元 $X$ 必须满足
\begin{equation*}
 0=\{X,\sigma_i\}
 =2a\sigma_i+2b_i\id(2)
 \qquad(i=1,2,3).
\end{equation*}
取迹给 $b_i=0$；再乘任意 $\sigma_j$ 后取迹给 $a=0$。所以唯一同时与三个 Pauli 方向反对易的二维矩阵是零矩阵，这与 $(\gamma^0)^2=\id(2)$ 矛盾。

因此非零复表示的维数既不能为奇数，也不能为二，最小可能值至少为四。正文\eqnrefs{eq:weyl-gamma} 又给出显式 $4\times4$ 矩阵并满足\eqnrefs{eq:clifford-general}，所以
\begin{equation*}
 \dim_{\mathbb C}S_{\rm Dirac}^{\rm min}=4.
\end{equation*}
这完成“最小维数为四”的存在性与不可更小性两部分证明。进一步，抽象 Clifford 关系可将任意单项式约化为四个不同生成元按序出现的乘积，故其维数至多为 $2^4=16$。上述四维表示中，这些乘积等价于已经证明完备的 Dirac 基
\[
I,\quad\gamma^\mu,\quad\sigma^{\mu\nu}\ (\mu<\nu),\quad\gamma^\mu\gamma^5,\quad\gamma^5.
\]
它们的迹配对非退化，因而是 $M_4(\CC)$ 的 $16$ 个线性无关矩阵。表示映射于是既满又单射，给出 $\mathrm{Cl}_4(\CC)\simeq M_4(\CC)$。最后，矩阵单位 $E_{ij}$ 将任一不可约模中非零的 $E_{jj}$ 分量生成一份标准模；不可约性使其覆盖整个模，故维数恰为四。偶子代数分解为 $M_2(\CC)\oplus M_2(\CC)$，对应两个 Weyl 模。
\end{proof}

\begin{proposition}[旋量 Lorentz 生成元的代数]
由 $\Sigma^{\mu\nu}=\ii[\gamma^\mu,\gamma^\nu]/4$ 定义的矩阵满足\eqnrefs{eq:spinor-lorentz-algebra}，并通过与 gamma 的交换子实现向量表示。
\end{proposition}

\begin{proof}
先计算
\begin{align*}
 [\gamma^\mu\gamma^\nu,\gamma^\rho]
 &=\gamma^\mu\gamma^\nu\gamma^\rho
 -\gamma^\rho\gamma^\mu\gamma^\nu\\
 &=\gamma^\mu(2\eta^{\nu\rho}-\gamma^\rho\gamma^\nu)
 -(2\eta^{\mu\rho}-\gamma^\mu\gamma^\rho)\gamma^\nu\\
 &=2\eta^{\nu\rho}\gamma^\mu-2\eta^{\mu\rho}\gamma^\nu.
\end{align*}
交换 $\mu\leftrightarrow\nu$ 后相减，再乘正文\eqnrefs{eq:sigma-definition} 中的 $\ii/4$，便得到
\begin{equation*}
 [\Sigma^{\mu\nu},\gamma^\rho]
 =\ii(\eta^{\nu\rho}\gamma^\mu-\eta^{\mu\rho}\gamma^\nu),
\end{equation*}
即正文\eqnrefs{eq:sigma-gamma-comm}。

由 $\Sigma^{\rho\sigma}=\frac{\ii}{4}[\gamma^\rho,\gamma^\sigma]$，有
\begin{align*}
 [\Sigma^{\mu\nu},\Sigma^{\rho\sigma}]
 &=\frac{\ii}{4}\Bigl(
 [\Sigma^{\mu\nu},\gamma^\rho]\gamma^\sigma
 +\gamma^\rho[\Sigma^{\mu\nu},\gamma^\sigma]\\
 &\qquad-[\Sigma^{\mu\nu},\gamma^\sigma]\gamma^\rho
 -\gamma^\sigma[\Sigma^{\mu\nu},\gamma^\rho]\Bigr).
\end{align*}
代入正文\eqnrefs{eq:sigma-gamma-comm}，再把出现的 $\gamma^\alpha\gamma^\beta-\gamma^\beta\gamma^\alpha$ 重新写回 $\Sigma^{\alpha\beta}$，得到
\begin{equation*}
 [\Sigma^{\mu\nu},\Sigma^{\rho\sigma}]
 =\ii\left(
 \eta^{\nu\rho}\Sigma^{\mu\sigma}
 -\eta^{\mu\rho}\Sigma^{\nu\sigma}
 -\eta^{\nu\sigma}\Sigma^{\mu\rho}
 +\eta^{\mu\sigma}\Sigma^{\nu\rho}
 \right),
\end{equation*}
即正文\eqnrefs{eq:spinor-lorentz-algebra}。
\end{proof}

\begin{proposition}[gamma 矩阵的有限协变律]
若 $S$ 是连通保向保时 Lorentz 变换 $\Lambda$ 的旋量提升，则 $S^{-1}\gamma^\rho S=\Lambda^\rho{}_\sigma\gamma^\sigma$。两个提升相差整体负号，均给相同相似变换。
\end{proposition}

\begin{proof}
正文\eqnrefs{eq:spinor-lorentz-transform} 给出有限旋量变换的指数形式。其无穷小相似变换由 $[\Sigma^{\mu\nu},\gamma^\rho]$ 完全决定；正文\eqnrefs{eq:sigma-gamma-comm} 恰好把这一对易子写成四矢量表示的生成元作用，因此两种表示具有同一 Lie 代数微分，指数化后得到 $\gamma^\rho$ 的有限协变变换律。

写成一阶参数展开，
\begin{equation*}
 S=\id()-\frac{\ii}{2}\omega_{\mu\nu}\Sigma^{\mu\nu}+O(\omega^2),
 \qquad
 S^{-1}=\id()+\frac{\ii}{2}\omega_{\mu\nu}\Sigma^{\mu\nu}+O(\omega^2).
\end{equation*}
于是
\begin{align*}
 S^{-1}\gamma^\rho S
 &=\gamma^\rho+\frac{\ii}{2}\omega_{\mu\nu}
 [\Sigma^{\mu\nu},\gamma^\rho]+O(\omega^2)\\
 &=\gamma^\rho+\omega^\rho{}_{\sigma}\gamma^\sigma+O(\omega^2).
\end{align*}
四矢量表示同样满足
\begin{equation*}
 \Lambda^\rho{}_{\sigma}
 =\delta^\rho{}_{\sigma}+\omega^\rho{}_{\sigma}+O(\omega^2).
\end{equation*}
两种表示由同一 Lie 代数参数生成，因此沿群参数连续积分后得到
\begin{equation*}
 S^{-1}(\Lambda)\gamma^\rho S(\Lambda)
 =\Lambda^\rho{}_{\sigma}\gamma^\sigma,
\end{equation*}
即正文\eqnrefs{eq:gamma-covariance}。
\end{proof}

\begin{proposition}[Dirac 共轭的协变律]
采用 $(\gamma^\mu)^\dagger=\gamma^0\gamma^\mu\gamma^0$，旋量提升满足 $S^\dagger\gamma^0=\gamma^0S^{-1}$。因此 $\bar\psi=\psi^\dagger\gamma^0$ 在 $\psi^\prime=S\psi$ 下变为 $\bar\psi^\prime=\bar\psi S^{-1}$。
\end{proposition}

\begin{proof}
普通内积 $\psi^\dagger\psi$ 在推动下并不保持不变，因此正文\eqnrefs{eq:dirac-adjoint} 引入 $\gamma^0$ 来构造适合 Lorentz 群的双线性。要验证这个定义确实有效，只需证明 $S^\dagger\gamma^0=\gamma^0S^{-1}$。

在所采用的度规约定下，
\begin{equation*}
 (\gamma^0)^\dagger=\gamma^0,
 \qquad
 (\gamma^i)^\dagger=-\gamma^i,
\end{equation*}
等价于
\begin{equation*}
 (\gamma^\mu)^\dagger=\gamma^0\gamma^\mu\gamma^0.
\end{equation*}
因此
\begin{equation*}
 (\Sigma^{\mu\nu})^\dagger=\gamma^0\Sigma^{\mu\nu}\gamma^0.
\end{equation*}
对正文\eqnrefs{eq:spinor-lorentz-transform} 取厄米共轭便有
\begin{align*}
 S^\dagger
 &=\exp\!\left[+\frac{\ii}{2}\omega_{\mu\nu}(\Sigma^{\mu\nu})^\dagger\right]\\
 &=\gamma^0\exp\!\left[+\frac{\ii}{2}\omega_{\mu\nu}\Sigma^{\mu\nu}\right]\gamma^0
 =\gamma^0S^{-1}\gamma^0.
\end{align*}
于是
\begin{align*}
 \bar\psi'
 &=(S\psi)^\dagger\gamma^0
 =\psi^\dagger S^\dagger\gamma^0
 =\psi^\dagger\gamma^0S^{-1}
 =\bar\psi S^{-1},
\end{align*}
这就是正文\eqnrefs{eq:dirac-adjoint-covariance-dp68}。
\end{proof}

\begin{lemma}[手征投影的代数]
四维中定义 $\gamma^5=\ii\gamma^0\gamma^1\gamma^2\gamma^3$，则 $(\gamma^5)^2=I$、$\{\gamma^5,\gamma^\mu\}=0$。$P_{R,L}=(I\pm\gamma^5)/2$ 为互补投影，且 $P_R\gamma^\mu=\gamma^\mu P_L$。
\end{lemma}

\begin{proof}
以 $\mu=0$ 为例，
\begin{align*}
 \gamma^5\gamma^0
 &=\ii\gamma^0\gamma^1\gamma^2\gamma^3\gamma^0\\
 &=-\ii\gamma^0\gamma^0\gamma^1\gamma^2\gamma^3
 =-\gamma^0\gamma^5.
\end{align*}
同理对任意 $\mu$ 都有 $\{\gamma^5,\gamma^\mu\}=0$。平方时，右半组四个 $\gamma$ 移到左半组同名矩阵旁共需 $3+2+1=6$ 次交换：
\begin{align*}
 (\gamma^5)^2
 &=-\gamma^0\gamma^1\gamma^2\gamma^3
 \gamma^0\gamma^1\gamma^2\gamma^3\\
 &=-(-1)^6(\gamma^0)^2(\gamma^1)^2(\gamma^2)^2(\gamma^3)^2
 =\id(4).
\end{align*}
两条关系合起来就是正文\eqnrefs{eq:gamma5-anticomm}。

先计算
\begin{align*}
 P_R^2&=\frac14(1+2\gamma^5+(\gamma^5)^2)=P_R,\\
 P_L^2&=\frac14(1-2\gamma^5+(\gamma^5)^2)=P_L,\\
 P_RP_L&=\frac14(1-(\gamma^5)^2)=0,\\
 P_R+P_L&=\id(4).
\end{align*}
再用 $\gamma^5\gamma^\mu=-\gamma^\mu\gamma^5$，
\begin{align*}
 P_R\gamma^\mu
 &=\frac12(1+\gamma^5)\gamma^\mu
 =\frac12\gamma^\mu(1-\gamma^5)
 =\gamma^\mu P_L,\\
 P_L\gamma^\mu&=\gamma^\mu P_R.
\end{align*}
这就得到正文\eqnrefs{eq:chiral-projector-algebra-dp68}。
\end{proof}

\begin{proposition}[幺正表示的 Hermitian 生成元闭合]
在光滑幺正群表示的共同不变光滑向量域上，按 $U(\theta)=I-\ii\theta^aT_a/\hbar+O(\theta^2)$ 取生成元，则 $[T_a,T_b]=\ii\hbar f_{ab}{}^cT_c$。有限维时 $T_a$ 为 Hermitian 矩阵；无界情形需以 Stone 定理定义自伴生成元，不能仅由形式伴随断言。
\end{proposition}

\begin{proof}
在量子态 Hilbert 空间中，把单位元附近的连续变换写成

\begin{equation*}
 U(\boldsymbol\theta)
 =\id()-\frac{\ii}{\hbar}\theta^aT_a+
 \order(\theta^2).
\end{equation*}
若 $U$ 对实参数保持内积，$U^\dagger U=\id()$ 的一阶项给
\begin{equation*}
 T_a^\dagger=T_a,
\end{equation*}
所以这里采用厄米生成元约定。

再取
\begin{equation*}
 U_a(\epsilon)=\ee^{-\ii\epsilon T_a/\hbar},
 \qquad
 U_b(\delta)=\ee^{-\ii\delta T_b/\hbar}.
\end{equation*}
正文\eqnrefs{eq:group-commutator-bch-dp8} 已由 BCH 公式给出
\begin{equation*}
 U_a(\epsilon)U_b(\delta)U_a(-\epsilon)U_b(-\delta)
 =\id()-\frac{\epsilon\delta}{\hbar^2}[T_a,T_b]+
 \order(\epsilon^2\delta,\epsilon\delta^2).
\end{equation*}
另一方面，这个乘积仍是群元素，而且在 $\epsilon,\delta\to0$ 时趋近单位元，所以必可写成某个新的无穷小参数 $\delta\theta^c=\order(\epsilon\delta)$：
\begin{equation*}
 U(\delta\boldsymbol\theta)
 =\id()-\frac{\ii}{\hbar}\delta\theta^cT_c+
 \order\!\left(\|\delta\boldsymbol\theta\|^2\right).
\end{equation*}
局部光滑性允许把最低阶参数写成
\begin{equation*}
 \delta\theta^c=\epsilon\delta\,f_{ab}{}^c,
\end{equation*}
其中 $f_{ab}{}^c$ 与 $\epsilon,\delta$ 无关。比较两个 $\epsilon\delta$ 阶系数便有
\begin{equation*}
 [T_a,T_b]=\ii\hbar f_{ab}{}^cT_c,
\end{equation*}
即正文\eqnrefs{eq:lie-algebra-closure-dp68}。结构常数的反对称性来自对易子的反对称性；Jacobi 恒等式则直接继承自算符对易子的 Jacobi 恒等式。
\end{proof}
